\documentclass[conference]{IEEEtran}
\IEEEoverridecommandlockouts

\usepackage[spanish,es-nodecimaldot,es-noshorthands]{babel}
\usepackage[utf8]{inputenc}
\usepackage[T1]{fontenc}
\usepackage{amsmath,amssymb}
\usepackage{graphicx}
\usepackage{xcolor}
\usepackage{float}
\usepackage{array}
\usepackage{listings}
\usepackage[hidelinks]{hyperref}

\definecolor{gris}{rgb}{0.45,0.45,0.45}
\definecolor{pend}{rgb}{0.75,0.10,0.10}

\lstset{
	basicstyle=\ttfamily\scriptsize,
	frame=single,
	breaklines=true,
	showstringspaces=false,
	columns=fullflexible,
	literate=%
	{á}{{\'a}}1 {é}{{\'e}}1 {í}{{\'i}}1 {ó}{{\'o}}1 {ú}{{\'u}}1
	{Á}{{\'A}}1 {É}{{\'E}}1 {Í}{{\'I}}1 {Ó}{{\'O}}1 {Ú}{{\'U}}1
	{ñ}{{\~n}}1 {Ñ}{{\~N}}1 {¿}{{?`}}1 {¡}{{!`}}1
}

% Recuadro para una imagen que todavía no existe: \imgph{descripción}{alto}
\newcommand{\imgph}[2]{%
	\fbox{\begin{minipage}[c][#2][c]{0.92\columnwidth}
			\centering\textcolor{gris}{\footnotesize [Espacio para imagen: #1]}
	\end{minipage}}%
}
% Marca en rojo lo que falta llenar. Antes de entregar no debe quedar ninguno.
\newcommand{\pend}[1]{\textcolor{pend}{[#1]}}

\begin{document}

\title{Avance 3: Integración de Extremo a Extremo \\
	\large Terminal de Juegos Retro Distribuida con Mando Embebido}

\author{
	\IEEEauthorblockN{Belmont Muñoz Samuel}
	\IEEEauthorblockA{No. Cuenta: 317318289}
	\and
	\IEEEauthorblockN{Domínguez Palacios Jesús Alejandro}
	\IEEEauthorblockA{No. Cuenta: 421066205}
	\and
	\IEEEauthorblockN{Dorantes Mejía Emmanuel}
	\IEEEauthorblockA{No. Cuenta: 313004766}
	\and
	\IEEEauthorblockN{García Olvera Emiliano}
	\IEEEauthorblockA{No. Cuenta: 314256009}
	\and
	\IEEEauthorblockN{Hernández Alonso Antonio Adonahí}
	\IEEEauthorblockA{No. Cuenta: 320257780}
	\and
	\IEEEauthorblockN{Jiménez Reyes Alan Yael}
	\IEEEauthorblockA{No. Cuenta: 320026726}
}

\maketitle

\begin{center}
	\vspace{-0.3cm}
	{\small \textbf{Fundamentos de Sistemas Embebidos, Grupo 6} \quad
		\textbf{Fecha:} 09/11/2026}
\end{center}

\begin{abstract}
Este avance integra todos los bloques del sistema: el mando basado en Raspberry Pi Pico lee 10 botones
mediante dos registros 74HC165 en cascada y un potenciómetro, y se presenta ante la Raspberry Pi 4 como un
control USB HID. La consola ejecuta RetroArch sobre Raspberry Pi OS Lite y muestra el juego en una pantalla
TFT SPI de 3.5''. Un proceso de telemetría lee la RAM del juego y, por medio de un socket UNIX y un segundo
proceso, ordena a la Pico activar el motor de vibración, el buzzer y los LEDs. Todo el sistema arranca solo
al encenderse mediante servicios de \texttt{systemd}.
\end{abstract}

\begin{IEEEkeywords}
Raspberry Pi Pico, Raspberry Pi 4, USB HID, RetroArch, systemd, socket UNIX, PIO, DMA.
\end{IEEEkeywords}

%=====================================================================
% SECCIÓN DE JESÚS: sistema integrado y demostración
%=====================================================================
\section{Sistema integrado}
\textbf{Autor:} Domínguez Palacios Jesús Alejandro

La Fig.~\ref{fig:bloques} muestra el diagrama de bloques actualizado. Todos los bloques del Avance 1 están
presentes y conectados; la Tabla~\ref{tab:bloques} indica cómo se implementó cada uno y en qué archivo está.

\begin{figure*}[!t]
	\centering
	\imgph{diagrama de bloques del Avance 3 (diagrama.py de Adonahí)}{6cm}
	% Cuando exista: \includegraphics[width=\textwidth]{../img/diagrama_bloques_v3.png}
	\caption{Diagrama de bloques del sistema integrado.}
	\label{fig:bloques}
\end{figure*}

\begin{table*}[!t]
	\centering
	\caption{Bloques del sistema y su implementación}
	\label{tab:bloques}
	\small
	\begin{tabular}{|p{3.2cm}|p{6.6cm}|p{4.6cm}|p{1.4cm}|}
		\hline
		\textbf{Bloque} & \textbf{Implementación} & \textbf{Archivo} & \textbf{Estado} \\
		\hline
		Botones y 74HC165 $\times 2$ & 10 botones en cascada, lectura por PIO + DMA cada 5 ms & \texttt{firmware/botones.c}, \texttt{read\_74hc165.pio} & \pend{ok?} \\
		\hline
		Potenciómetro & ADC de 12 bits, cada 10 ms & \texttt{firmware/adc\_pot.c} & \pend{ok?} \\
		\hline
		Enlace USB & Dispositivo compuesto: HID gamepad + CDC serial (TinyUSB) & \texttt{firmware/main.c}, \texttt{usb\_descriptors.c} & \pend{ok?} \\
		\hline
		Actuadores & Motor y buzzer por PWM, LEDs por GPIO; órdenes \texttt{VIB}, \texttt{BUZ}, \texttt{LED} & \texttt{firmware/actuadores.c} & \pend{ok?} \\
		\hline
		Consola (Linux) & Raspberry Pi OS Lite, kernel 6.18, arranque con \texttt{systemd} & \texttt{pi/systemd/*.service} & \pend{ok?} \\
		\hline
		Emulación & RetroArch con driver \texttt{gl} sobre DRM & \texttt{pantalla/ra\_gl.cfg} & \pend{ok?} \\
		\hline
		Pantalla TFT 3.5'' & ST7796S por SPI a 50 MHz, driver \texttt{panel-mipi-dbi}, 480$\times$320 & \texttt{pantalla/config\_txt\_pantalla.txt} & Lista \\
		\hline
		Telemetría e IPC & Lectura de RAM por UDP 55355; socket \texttt{AF\_UNIX} hacia el puente serial & \texttt{pi/telemetria.py}, \texttt{pi/puente\_serial.py} & \pend{ok?} \\
		\hline
		Red LAN & Ethernet de la Pi 4; netplay en el Avance 4 & --- & Parcial \\
		\hline
	\end{tabular}
\end{table*}

\subsection{Flujo de datos de extremo a extremo}
\begin{enumerate}
	\item \textbf{Entrada:} un botón cambia de nivel; la alarma de 5 ms dispara la PIO, que serializa los 16 bits
	de los 74HC165; el DMA los copia a SRAM y la estructura \texttt{estado} se actualiza bajo un spin lock.
	\item \textbf{Hacia la consola:} Core 0 arma el reporte HID y lo envía por USB; en Linux, el driver
	\texttt{usbhid} lo publica en \texttt{/dev/input/event}\pend{N}, de donde lo lee RetroArch.
	\item \textbf{Salida de video:} RetroArch dibuja el cuadro y el driver \texttt{panel-mipi-dbi} lo envía por SPI a la pantalla.
	\item \textbf{Retroalimentación:} \texttt{telemetria.py} lee las vidas y el puntaje; al detectar un cambio manda una
	orden por el socket UNIX a \texttt{puente\_serial.py}, que la escribe en \texttt{/dev/ttyACM0}; la Pico activa el actuador.
\end{enumerate}

\subsection{Demostración}
\begin{figure}[H]
	\centering
	\imgph{foto del sistema completo: control, Pi 4 y pantalla con el juego}{4.5cm}
	\caption{Sistema integrado durante la demostración.}
	\label{fig:demo}
\end{figure}

\begin{figure}[H]
	\centering
	\imgph{captura de \texttt{evtest} con botones y potenciómetro}{4cm}
	\caption{Eventos del control recibidos por Linux.}
	\label{fig:evtest}
\end{figure}

\subsection{Configuración reproducible}
Los pasos completos están en el \texttt{README.md} del repositorio. En resumen, en una Pi recién instalada:
\begin{lstlisting}
git clone https://github.com/YizusDP/Practicas_Embebidos.git
cd Practicas_Embebidos/ProyectoFinal
sudo ./pi/instalar.sh
sudo reboot
\end{lstlisting}
Tras el reinicio, \texttt{systemctl status arcade.target} muestra los tres servicios activos (Fig.~\ref{fig:arranque}).
\begin{figure}[H]
	\centering
	\imgph{captura de \texttt{systemctl status arcade.target} justo después de encender}{3.5cm}
	\caption{Arranque automático del sistema.}
	\label{fig:arranque}
\end{figure}

%=====================================================================
% SECCIONES DE CADA INTEGRANTE (cada quien edita solo la suya)
%=====================================================================
%=====================================================================
% Sección de Belmont Muñoz Samuel: pruebas de integración
% Los valores MEDIDOS (en rojo) los llena Samuel con el osciloscopio o el
% analizador lógico. No inventar datos.
%=====================================================================
\section{Pruebas de integración}
\textbf{Autor:} Belmont Muñoz Samuel

\textbf{Requisito medido:} RF1 del Avance 1, lectura de botones con periodo de $5\text{ ms} \pm 0.5\text{ ms}$ ($200\text{ Hz}$).

\textbf{Cálculo teórico.} La lectura la hace la PIO sin intervención de la CPU, así que el periodo lo fija
el número de ciclos del programa \texttt{read\_74hc165.pio} y el divisor de reloj
(\texttt{BTN\_SM\_CLKDIV = 125}, o sea $125\text{ MHz}/125 = 1\text{ MHz}$, $1\,\mu\text{s}$ por ciclo):
\begin{align*}
	\text{carga y liberación} &= 4 \times 32 = 128 \\
	\text{primer bit} + \texttt{set x} &= 1 + 1 = 2 \\
	\text{9 bits restantes} &= 9 \times (4 \times 32 + 1 + 1) = 1170 \\
	\text{pausa} &= 1 + 2 \times (1 + 32 \times 33 + 1) = 2117 \\
	\text{total} &= 3417 \text{ ciclos} = 3.417\text{ ms}\ (292.7\text{ Hz})
\end{align*}

\textbf{Método:} la PIO ejecuta exactamente una instrucción por ciclo de su reloj, sin interrupciones ni caché
que la retrasen, así que el periodo de la trama es determinista y queda fijado por el conteo de ciclos y el divisor.
La única fuente de variación es el oscilador de cristal de 12 MHz de la Pico, del orden de partes por millón.
La medición con osciloscopio en GPIO17 (señal LD) queda como confirmación para el Avance 4.

\begin{table}[H]
	\centering
	\caption{Periodo de lectura de botones contra el Avance 1}
	\label{tab:medicion}
	\small
	\begin{tabular}{|l|c|c|c|}
		\hline
		\textbf{Dato} & \textbf{Avance 1} & \textbf{Actual} & \textbf{Ajustado} \\
		\hline
		Periodo & $5.0 \pm 0.5\text{ ms}$ & $3.417\text{ ms}$ & $5.000\text{ ms}$ \\
		\hline
		Frecuencia & $200\text{ Hz}$ & $292.7\text{ Hz}$ & $200.0\text{ Hz}$ \\
		\hline
	\end{tabular}
\end{table}

\textbf{Resultado:} con el divisor actual (\texttt{BTN\_SM\_CLKDIV = 125}) la trama dura 3.417 ms, fuera del
intervalo de $5 \pm 0.5$ ms comprometido: \textbf{no cumple}, por ser 1.583 ms más rápida. Leer más rápido no afecta
la latencia, pero el requisito pide un periodo fijo. La corrección es cambiar el divisor a 182.91:
\[
T = \frac{3417 \times 182.91}{125\text{ MHz}} = 5.000\text{ ms}
\]
El divisor fraccionario de la PIO tiene resolución de $1/256$, así que $182.91 \approx 182 + 233/256 = 182.910$,
y el error queda por debajo del 0.01\,\% (menos de $0.5\,\mu\text{s}$).
          % Pruebas de integración: medición
%=====================================================================
% Sección de Hernández Alonso Antonio Adonahí: robustez (2 fallas inducidas)
% Solo Adonahí edita este archivo. Jesús ejecuta las fallas en la Pi y le pasa las capturas.
%=====================================================================
\section{Robustez: fallas inducidas}
\textbf{Autor:} Hernández Alonso Antonio Adonahí

\begin{table}[H]
	\centering
	\caption{Fallas inducidas y recuperación}
	\label{tab:fallas}
	\small
	\begin{tabular}{|p{1.7cm}|p{2.2cm}|p{2.4cm}|p{1.0cm}|}
		\hline
		\textbf{Falla} & \textbf{Cómo se induce} & \textbf{Cómo se recupera} & \textbf{Tiempo} \\
		\hline
		F1: proceso muerto & \texttt{sudo kill -9 \$(pidof retroarch)} & \texttt{systemd} lo relanza por \texttt{Restart=always} & \pend{s} \\
		\hline
		F2: control desconectado & Se desconecta el USB del control en plena partida & RetroArch lo reconoce al reconectarlo; el puente reabre \texttt{/dev/ttyACM0} & \pend{s} \\
		\hline
	\end{tabular}
\end{table}

\begin{figure}[H]
	\centering
	\imgph{captura de \texttt{journalctl -u arcade-retroarch} con la caída y el reinicio (F1)}{3.5cm}
	\caption{Recuperación de la falla F1.}
	\label{fig:f1}
\end{figure}

\begin{figure}[H]
	\centering
	\imgph{captura de \texttt{dmesg} o del log del puente al reconectar el control (F2)}{3.5cm}
	\caption{Recuperación de la falla F2.}
	\label{fig:f2}
\end{figure}
  % Robustez: 2 fallas inducidas

\section{Bitácora: planeado frente a construido}
La Tabla~\ref{tab:bitacora} resume dónde quedó cada decisión. Las entradas firmadas de cada clase están a continuación.

\begin{table*}[!t]
	\centering
	\caption{Decisiones de las clases 8, 9 y 11: planeado frente a construido}
	\label{tab:bitacora}
	\small
	\begin{tabular}{|p{1.2cm}|p{3.6cm}|p{3.8cm}|p{3.2cm}|p{3.8cm}|}
		\hline
		\textbf{Clase} & \textbf{Planeado (Avance 1)} & \textbf{Construido} & \textbf{Archivo:línea} & \textbf{Diferencia y causa} \\
		\hline
		8 Pico & Lectura de botones por interrupción de temporizador cada 5 ms, prioridad alta & \pend{Samuel} & \pend{botones.c:N} & \pend{} \\
		\hline
		8 Pi 4 & RetroArch con \textit{nice} menor que la telemetría & \texttt{Nice=-5} en RetroArch, \texttt{Nice=10} en telemetría & \pend{arcade-retroarch.service:N} & Ninguna \\
		\hline
		9 Pico & Mutex de hardware sobre el reporte HID & Spin lock sobre \texttt{estado} & \pend{estado.h:N, botones.c:N} & No se puede dormir esperando un mutex dentro de una interrupción \\
		\hline
		9 Pi 4 & Triple buffer con semáforos POSIX & Lo resuelve RetroArch con DRM & \texttt{ra\_gl.cfg:3261} & No es código nuestro; se configura con \texttt{video\_max\_swapchain\_images = "3"} \\
		\hline
		11 & Socket \texttt{AF\_UNIX DGRAM} expuesto por RetroArch & RetroArch expone UDP 55355; el socket UNIX une telemetría y puente & \pend{telemetria.py:N, puente\_serial.py:N} & RetroArch no ofrece socket UNIX \\
		\hline
	\end{tabular}
\end{table*}

%=====================================================================
% Sección de Hernández Alonso Antonio Adonahí: bitácora clases 8 y 9
%=====================================================================
\subsection{Clase 8: Planificación}
\textbf{Autor:} Hernández Alonso Antonio Adonahí

\textbf{Planeado:} en el Avance 1 se definieron tres tareas principales con sus prioridades: la emulación (RetroArch) con la máxima prioridad para garantizar 60 FPS constantes, la adquisición de botones a 200 Hz (5 ms) con prioridad media para cumplir el requisito de latencia, y la telemetría con prioridad baja por no ser crítica para la jugabilidad.

\textbf{Construido:}
\begin{itemize}
	\item \textbf{Pico:} la alarma de 5 ms que dispara la lectura de los registros 74HC165 se encuentra en \texttt{firmware/botones.c}.
	\item \textbf{Pi 4:} \texttt{Nice=-5} para RetroArch (\texttt{pi/systemd/arcade-retroarch.service}, línea 18) y \texttt{Nice=10} para la telemetría (\texttt{pi/systemd/arcade-telemetria.service}, línea 11).
\end{itemize}

\begin{figure}[H]
	\centering
	\includegraphics[width=0.85\columnwidth]{avance3/img/arranque_automatico.png}
	\caption{Prioridades de los procesos en la consola mostrando RetroArch con NI -5 y telemetría con NI 10.}
	\label{fig:nice}
\end{figure}

\textbf{Diferencias:} el \texttt{Nice} de la telemetría es 10 en lugar de 5 como se planeó originalmente, para asegurar que nunca interfiera con la emulación.

\subsection{Clase 9: Sincronización}
\textbf{Autor:} Hernández Alonso Antonio Adonahí

\textbf{Planeado:} mutex de hardware en la Pico para proteger la estructura compartida entre el Core 0 (USB) y la alarma que lee los botones; triple buffer con semáforos POSIX en la Pi para la comunicación entre procesos.

\textbf{Construido:}
\begin{itemize}
	\item \textbf{Pico:} no se requirió spin lock ni mutex. El DMA escribe cada campo de \texttt{g\_controls} con una sola transferencia de 16 bits y el USB copia campo por campo con lecturas atómicas de 16 bits (\texttt{controls\_snapshot} en \texttt{usb\_data\_tx.c}), garantizando que ningún valor se lea a medias.
	\item \textbf{Pi:} el puente es el único proceso que escribe en el dispositivo HID del control (\texttt{/dev/arcade\_hid}), eliminando carreras de escritura por diseño. El triple buffer lo implementa RetroArch mediante DRM (\texttt{ra\_gl.cfg}, línea 3261, \texttt{video\_max\_swapchain\_images = "3"}).
\end{itemize}

\textbf{Diferencias:} se planeó usar mutex de hardware en la Pico, pero el diseño final con transferencias atómicas de 16 bits eliminó la necesidad de sincronización explícita. Los semáforos POSIX planeados para el triple buffer no fueron necesarios porque RetroArch ya implementa su propio triple buffer a nivel de driver gráfico.
   % Clases 8 y 9
%===========================================================
% Sección de Jiménez Reyes Alan Yael: bitácora clase 11 (IPC)
%===========================================================
\subsection{Clase 11: Comunicación entre procesos}
\textbf{Autor:} Jiménez Reyes Alan Yael

\textbf{Planeado:} RetroArch expone \texttt{/tmp/retroarch\_telemetry.sock} (\texttt{AF\_UNIX}, \texttt{SOCK\_DGRAM}) y un script de telemetría se conecta a él.

\textbf{Construido:}
\begin{itemize}
    \item \texttt{telemetria.py} pide la RAM del juego a RetroArch por UDP en \texttt{127.0.0.1:55355} (línea 50).
    \item Envía cada orden como datagrama a \texttt{/tmp/arcade\_act.sock} (línea 61).
    \item \texttt{puente\_hid.py} recibe el datagrama (socket en la línea 73), lo traduce a un reporte de salida HID de 8 bytes (línea 46) y lo escribe en \texttt{/dev/arcade\_hid} (línea 105)[cite: 23].
\end{itemize}

\textbf{Direcciones de RAM usadas:} Pac-Man (core Nestopia), dirección \texttt{0x067}, significa el contador de vidas del jugador (empieza en 3 y baja 1 con cada muerte).

\textbf{Diferencias:} RetroArch no ofrece un socket UNIX, sino un servidor de comandos por UDP en el puerto 55355; por lo tanto, el socket UNIX quedó implementado de manera interna como un puente entre nuestros dos procesos de Python. En el Avance 3 primero se usó un serial CDC para las órdenes; se cambió a reportes de salida HID para respetar el diagrama (un solo enlace USB HID)[cite: 23].

\subsection{Declaración de uso de IA}
Para la elaboración de este avance, utilicé asistencia de Inteligencia Artificial para verificar la correcta ejecución de los commits solicitados y estructurar los textos de esta bitácora.
      % Clase 11
%=====================================================================
% Sección de García Olvera Emiliano: bitácora clase 20 (tipos y volatile)
%=====================================================================
\subsection{Clase 20: Tipos y \texttt{volatile}}
\textbf{Autor:} García Olvera Emiliano

\textbf{Pregunta:} \pend{copiar la pregunta de la clase 20}.

\textbf{Respuesta aplicada al proyecto:}
\begin{itemize}
	\item \texttt{uint16\_t botones}: \pend{por qué 16 bits}.
	\item \texttt{uint16\_t pot}: \pend{12 bits del ADC caben en 16}.
	\item \texttt{sizeof(estado\_t)} $=$ \pend{N} bytes; reporte HID de \pend{N} bytes.
	\item \texttt{volatile} en \pend{archivo, línea}: \pend{por qué}.
	\item Tamaño del firmware: \pend{Flash y RAM de \texttt{arm-none-eabi-size}}.
\end{itemize}

\begin{figure}[H]
	\centering
	\imgph{captura de \texttt{arm-none-eabi-size} y de \texttt{sizeof} impreso por la Pico}{3cm}
	\caption{Evidencia de tipos y tamaños.}
	\label{fig:tipos}
\end{figure}
  % Clase 20: tipos y volatile
%=====================================================================
% Sección de Dorantes Mejía Emmanuel: bitácora clase 21 (drivers)
%=====================================================================
\subsection{Clase 21: Driver de cada sensor y actuador}
\textbf{Autor:} Dorantes Mejía Emmanuel

\textbf{Pregunta:} Para cada sensor o actuador de su proyecto,
indiquen qué archivo de dispositivo o biblioteca usa y qué
driver del kernel hay detrás (con \texttt{ls -l /dev},
\texttt{lsmod} o \texttt{dmesg} como evidencia).
Si alguna parte necesitara un driver propio, expliquen
por qué no alcanza el genérico.

\medskip
\textbf{Respuesta:}

El proyecto utiliza una Raspberry Pi Pico para adquirir las
señales de entrada y controlar los actuadores, mientras que
la Raspberry Pi 4 ejecuta Linux y gestiona la comunicación
USB y la visualización del videojuego.

En la Pico, los periféricos se controlan mediante el SDK
y código propio, sin utilizar drivers del kernel Linux.
En la Raspberry Pi 4 se emplean los controladores
correspondientes a los dispositivos conectados.

\begin{table}[H]
    \centering
    \caption{Drivers de sensores y actuadores}
    \label{tab:drivers}
    \small
    \begin{tabular}{|p{1.8cm}|p{2.6cm}|p{2.6cm}|}
        \hline
        \textbf{Elemento} & \textbf{Driver} &
        \textbf{Evidencia} \\
        \hline
        Botones &
        PIO + DMA (propio) &
        \texttt{botones.c} y
        \texttt{read\_74hc165.pio} \\
        \hline
        Potenciómetro &
        \texttt{hardware/adc} y
        \texttt{hardware/dma} &
        \texttt{read\_adc.c} y
        \texttt{read\_adc.h} \\
        \hline
        Motor, buzzer &
        \texttt{hardware/pwm} &
        \texttt{actuadores.c} \\
        \hline
        LEDs &
        \texttt{hardware/gpio} &
        \texttt{actuadores.c} \\
        \hline
        Control (Pi) &
        \texttt{usbhid}
        $\rightarrow$
        \texttt{/dev/input/eventX} &
        \texttt{/proc/bus/input/devices}
        (pendiente de verificar) \\
        \hline
        Serial (Pi) &
        \texttt{cdc\_acm}
        $\rightarrow$
        \texttt{/dev/ttyACM0} &
        \texttt{dmesg}
        (pendiente de verificar) \\
        \hline
        Pantalla (Pi) &
        \texttt{panel-mipi-dbi} &
        \texttt{lsmod}
        (pendiente de verificar) \\
        \hline
    \end{tabular}
\end{table}

\medskip
\textbf{Aportación individual:}

Se desarrollaron los archivos \texttt{read\_adc.c} y
\texttt{read\_adc.h} para implementar la adquisición
continua de señales analógicas mediante el ADC de
12 bits de la Raspberry Pi Pico.

La implementación configura los canales ADC0, ADC1,
ADC2 y ADC3 mediante muestreo round-robin, con una
frecuencia configurada de 4000 muestras por segundo
por canal.

Se utilizan dos canales DMA encadenados para transferir
automáticamente las muestras hacia un búfer circular
de cuatro elementos de 16 bits, reduciendo la
intervención de la CPU.

El archivo \texttt{read\_adc.c} contiene la inicialización
del ADC, la configuración del DMA y la función
\texttt{read\_adc\_is\_alive()}, que permite supervisar
el avance de las transferencias.

Por su parte, \texttt{read\_adc.h} define las constantes
de configuración y declara las funciones utilizadas
por el módulo.

Esta implementación utiliza los periféricos ADC y DMA
del microcontrolador mediante las bibliotecas del
Pico SDK, sin requerir un driver del kernel Linux.

Para la lectura de botones se utiliza código propio
basado en PIO y DMA, debido a que se requiere controlar
la señal de carga paralela y el reloj de los registros
74HC165. Los demás periféricos emplean las bibliotecas
y controladores indicados en la tabla.

\begin{figure}[H]
    \centering
    \imgph{Capturas pendientes de
    \texttt{lsusb},
    \texttt{dmesg | tail},
    \texttt{cat /proc/bus/input/devices}
    y \texttt{lsmod | grep mipi}}{4cm}
    \caption{Verificación de los controladores
    de dispositivos en la Raspberry Pi 4.}
    \label{fig:drivers}
\end{figure}
  % Clase 21: drivers
%=====================================================================
% Sección de Jiménez Reyes Alan Yael: pendientes del Avance 4 y declaración de IA
%=====================================================================
\section{Avance contra el plan}
\textbf{Autor:} Jiménez Reyes Alan Yael

\subsection{Pendientes para el Avance 4}
\begin{itemize}
	\item \textbf{Netplay en LAN:} partida entre dos consolas; medir RNF3 (ping $<$ 5 ms y 0\,\% de pérdida).
	\item \textbf{Latencia de extremo a extremo:} medir del botón a la pantalla contra el requisito de $\leq$ 15 ms.
	\item \textbf{Periodo de lectura:} ajustar el divisor de la PIO para cumplir $5 \pm 0.5$ ms y volver a medir.
	\item \textbf{Segundo juego:} agregar Pong usando el potenciómetro como paleta.
	\item \textbf{Gabinete:} imprimir en 3D las piezas ya modeladas y montar el sistema.
	\item \textbf{Robustez:} raíz del sistema en solo lectura para soportar apagones sin corromper la microSD.
\end{itemize}

\subsection{Desviaciones del plan y su causa}
\begin{table}[H]
	\centering
	\caption{Desviaciones respecto al Avance 1}
	\label{tab:desviaciones}
	\small
	\begin{tabular}{|p{2.3cm}|p{5.2cm}|}
		\hline
		\textbf{Planeado} & \textbf{Construido y causa} \\
		\hline
		SPI para leer el 74HC165 & PIO + DMA: el SPI no genera el pulso de carga paralela que necesita el registro. \\
		\hline
		Interrupción de temporizador cada 5 ms & La PIO y el DMA leen solos, sin CPU. El periodo quedó en 3.42 ms por el divisor actual (ver Pruebas de integración). \\
		\hline
		Mutex sobre el reporte HID & Sin candado: el DMA escribe campos de 16 bits en una sola transferencia y un mutex no puede proteger contra el DMA. \\
		\hline
		RetroArch expone un socket UNIX & RetroArch expone comandos por UDP 55355; el socket UNIX quedó entre telemetría y puente. \\
		\hline
		Triple buffer con semáforos POSIX & Lo resuelve RetroArch con DRM (\texttt{ra\_gl.cfg}, línea 3261); no es código del equipo. \\
		\hline
		Órdenes a la Pico por serial CDC (primera versión del Avance 3) & Reporte de salida HID de 8 bytes, para respetar el diagrama con un solo enlace USB HID. \\
		\hline
		Antirrebote con red RC & $\tau = 330\,\Omega \times 100\text{ nF} = 33\,\mu\text{s}$, mucho menor que el rebote de un botón (1 a 10 ms); no sirve como antirrebote. \\
		\hline
		Pong y Pac-Man & Solo Pac-Man (NES, core Nestopia), el único core instalado; Pong pasa al Avance 4. \\
		\hline
		Audio del juego & Desactivado (\texttt{audio\_driver = "null"}): el prototipo no tiene bocina. \\
		\hline
		Netplay en el Avance 3 & Pasa al Avance 4 por tiempo. \\
		\hline
	\end{tabular}
\end{table}

\section{Declaración de uso de IA}
El equipo usó asistentes de IA generativa como apoyo durante este avance:
\begin{itemize}
	\item \textbf{Claude (Anthropic):} planeación y reparto de tareas, estructura del reporte en LaTeX,
	borradores de \texttt{telemetria.py}, \texttt{puente\_hid.py}, los servicios de \texttt{systemd},
	\texttt{instalar.sh}, el \texttt{README}, \texttt{actuadores.c} y el reporte de salida HID; revisión y
	compilación del firmware.
	\item \pend{Gemini (Google): ejemplo del formato del reporte de salida HID. Confirmar herramienta y uso.}
\end{itemize}
Todo el código generado se revisó y se probó en el hardware del equipo: RetroArch y la pantalla TFT,
la lectura de la RAM del juego (dirección \texttt{0x067}), el arranque automático con \texttt{systemd},
las prioridades (\texttt{Nice}) y la recuperación ante la falla F1 (3.4 s). El firmware compila sin
errores ni advertencias con el Pico SDK 2.1.1. Las decisiones de diseño, las mediciones y las respuestas
de la bitácora son responsabilidad de cada integrante.
\pend{Revisar que la lista refleje lo que realmente usó cada integrante.}
 % Pendientes del Avance 4 y declaración de IA

\vspace{0.2cm}
\noindent\textbf{Repositorio:} \\
\url{https://github.com/YizusDP/Practicas_Embebidos/tree/main/ProyectoFinal}

\end{document}
